\section{The idea of pre-existence}

To continue the translation of \textit{Classical and Romantic ethics}, we present the next two segments. In the first, \textbf{Julius Evola} describes how Catholicism recapitulates within itself the choice between Tradition and modernity. In the second, he provides a succinct summary of the Traditional path. Evola writes:

\begin{quotex}
The central view of Catholicism is that God, though creating man from nothing, allowed a miracle to happen, through which this being created from nothing is free, in the sense that he can rejoin himself to the root of his own being, to God, or else to deny the root, to constitute himself, to dissipate himself, to degenerate into a useless free will of a creature. This same doctrine, with due transpositions, can be applied to relations between the individual being and the spiritual entity, of which he is the creation and the human manifestation. We mean that the individual being, within given limits, uses equally his free will and that the same alternative is posed to him: either to will his own nature, to deepen it, and realize it until rejoining the pre-human and super-individual principle that corresponds to it; or else to give himself over to arbitrarily construct a mode of unnatural being, deprived of connections with those deepest forces or even in contradiction with them. This is exactly the existing opposition between the traditional, especially Nordic-Aryan, ideal and the modern ideal of civilization. For the former, the essential task is to know and to be oneself; for the latter, the task is instead to ``construct oneself'', to become what one is not, to shatter every limit in order to make everything possible to everyone: liberalism, democracy, individualism, ethical activist Protestantism, anti-racism, anti-traditionalism.
\end{quotex}


This is a fascinating and insightful passage. Here he shows that \textbf{Catholicism}, while expressing itself in theological rather than metaphysical terms, represents a tipping point\footnote{\url{http://en.wikipedia.org/wiki/The\_Tipping\_Point}}. The two possible understandings of that doctrine correspond precisely to the opposition between the Traditional and the modern worldview. That self-understanding could have tipped either way. Despite the efforts of popes Pius IX (Syllabus of Errors)\footnote{\url{http://www.papalencyclicals.net/Pius09/p9syll.htm}} and Pius X (On the Doctrine of the Modernists)\footnote{\url{http://www.papalencyclicals.net/Pius10/p10pasce.htm}} to limit the effects of modernism over a century ago, it seems that the trend now is to fully embrace modernism; that was confirmed with Vatican II which opened up to the modern world rather than continue its opposition. Nevertheless, it is still too slow for modern minds.

Here, Evola precisely represents the difference between the two world views. The Traditional man strives to actualize his True Self, his given nature. The modern man does not believe in such essences or formal causes, hence he does not have an essence. This is perfectly expressed by John Paul Sartre, one of the recent creators of the modern mind. Without an essence, he can be whatever he chooses. Hence, we see recent trends in the direction of gender-neutral child rearing\footnote{\url{http://www.independent.co.uk/news/world/europe/genderneutral-preschool-accused-of-mind-control-2305983.html}} or trans-gendered beauty contests\footnote{\url{http://www.mirror.co.uk/news/world-news/transgender-miss-universe-hopeful-jenna-786205}}. These trends are logical and rational within the modern worldview, so attempts to argue against them on the basis of their absurdity are futile. Even claiming that biological givens override human choice is useless. First of all, biological differences are regarded as unjust; the purpose of technology is to overcome such injustices. Besides, a biological defense of human nature has nothing at all to do with the Traditional understanding of it.

It is also important to note that the movements that Evola opposes are the very ones that most people today would define as the ultimate achievement of the West; hence, they define Western civilization to most minds. Evola continues:

\begin{quotex}
As it was traditionally taught, the doctrine of pre-existence therefore brings us beyond both fatalism as well as a poorly understood and individualistic freedom. Passing on to the more immediate consequences, in realizing his own nature, the individual harmonizes his human will with the super-human will that corresponds to it, he ``remembers'' himself, he establishes the connection with a principle that, being beyond birth, is likewise beyond death and every temporal condition: therefore, according to the ancient Indo-Aryan conception, this is the way by which, through action, he seeks to achieve ``liberation'' and to realize the divine. \emph{Dharma}---which signifies his own nature, duty, fidelity to blood, tradition, and caste---connects itself here, as we explained in our other book ``Revolt'' , to the sensation of being joined here from afar and does not signify limitations, as ``evolved spirits'' believe, but liberation. Led back to this traditional vision of life, all the principal causes of race acquire a higher and spiritual significance and the objection based on birth as either by chance or destiny loses every force.
\end{quotex}


There is little to add here. The idea of pre-existence cannot be understood as fatalism, while it stands in opposition to the modern idea of freedom from essences. When this is understood, man understands his True Will, he remembers who he is, which is the unmoved mover, the silent observer, beyond the contingent events of space and time. He achieves ultimate liberation which is the realization of the divine, the stage of \emph{theosis}. This connection with Dharma is not a limitation as believe the many moderns who consider themselves to be evolved.


\flrightit{Posted on 2012-04-11 by Cologero}

\begin{center}* * *\end{center}

\begin{footnotesize}
\begin{sffamily}
\texttt{Jason-Adam on 2013-07-12 at 18:22 said:}

Is it possible to imagine that, in the case of the transgender people you mentioned, that they are in fact expressing their True Self ? My question is can a person's Essence be of a differing nature from their body ? I read once on an anthroposophical website defending transgenderism using the reasoning that there is a spirit-body conflict. Perhaps part of the Kali Yuga running down of the world means people are placed in bodies they weren't supposed to be ?

\hfill

\texttt{Ash on 2013-07-13 at 03:53 said:}

Jason-Adam, I've read similar explanations, along the lines of the Evola-inspired analysis of same-sex relations being that the participants are in fact fulfilling the masculine-feminine duality with regards to their True Selves, which have not been expressed accurately in the body. It's an interesting point of view and one that would be worth exploring further, at least to go beyond the knee-jerk reactions (one way or another) all to common when these subjects come up. Certainly there are instances of such individuals being given places in past societies, including India with the Hijra, mentioned in the Kama Sutra.

\hfill

\texttt{Cologero on 2013-07-13 at 16:43 said:}

J-A and Ash. If you want to explore the issue further, there are some factors to consider as a necessary starting point. As to the so-called ``spirit body conflict'', the Anthroposophists are a unreliable starting point because of their extreme cartesianism. For them, the spirit inhabits the body much like a hermit crab looking for a shell. Also, for them, the various soul layers are independent of each other; hence their creepy idea of two Christ boys involving some sort of exchange of ``etheric bodies''. Our view is that the soul is the form of the body, hence they are intimately related. The rational soul, or spirit, may be transcendent to the lower souls and body, yet is still their ideal form.

In the Kali Yuga, there are possibilities of manifestation that are compossible at this point in the cycle, although that may not have been possible at other times. Nevertheless, it changes nothing about how we consider the ideal man, if you recall Evola's essay on involution and evolution. According to that view, beings may appear, but as deviations from the central axis. If they have their place, and they do since we reject the alternative, the divisive issue is to define precisely what that place should be. The modern world is unwilling to make any such distinctions about the distance from the ideal since it claims that all appearances are inherently equal and there is no ideal.

It is an effort to express one's ``true self'' which is always transcendent. It is glib to presume that each and every action is such an expression without considering it may actually be the expression of something less than the true self, maybe not even from the self at all but from exterior influences. In our time, some have chosen a difficult path for their ``true self'' but nevertheless that is the path to follow.

\hfill

\texttt{Jason-Adam on 2013-07-14 at 15:42 said:}

Cologero: We've stumbled upon something important now, the two ways of looking the relations of corporal and noncorporeal in the person, the one way of seeing them as coexisting the other as actually being related. I have personal experience with mind-body wars so Steiner's ideas resonated with me. Also the quote about Plotinus from Poryphory saying he felt as if his body was a prison feels true as well. Am I getting too close to Gnosticism though ?

As for the true self, it is indeed hard to tell what is Me and is not Me and what is inspired by evil unearthly beings. Meditation on myself has helped me, and I knew that I when reached a point of inner calm and strength based on what I envisioned that I had found myself and my life was forever given a new direction.

as for transgenders, what do you think of Cicero's claiming they had magic powers ? I've known some trans witches actually who seemed to have genuine abilities... ... ... .

\hfill

\texttt{Logres on 2013-07-14 at 21:49 said:}

Jason-Adam, I would imagine that this is true; I've seen a gay man argue on a news show that it isn't natural and isn't supposed to be idealized: such a person definitely has a certain nobleness, which many ``moral conservatives'' lack. He is sufficiently detached to intuit that it is ``wrong'', but not know why, but still condemn the idolizing and normalizing of it for youth.

\hfill

\texttt{Logres on 2013-07-14 at 21:51 said:}

Another thought: one's body could be viewed as a kind of ``oath'' the true Self took, rightly or wrongly. So what you are asking is, is it ever in some manner justified to violate an oath? I would imagine that there would be consequences.

\hfill

\texttt{Jason-Adam on 2013-07-15 at 21:36 said:}

Plato believed we all choose our bodies, I am not so sure. My feelings is the during Kali Yuga the process of birth becomes more mechanised and the personal equation (talking in the astral realm ) declines and more ``errors'' result than in previous days. Of course the ideal remains the same.

The Buddhists claim that transgender (kathoey) is a punishment for being an adulteress in the previous life but I dont believe in reincarnation.

Some traditional scholars such as Imam Khomeini (though is he traditional ?) support sex changes and oppose homosexuality.

As I understand it the purpose of sex and marriage is the unity of yin and yang, positive and negative. same sex attraction does not do this but what if the person is of the other gender in another body ?

I am Catholic and know the Church identifies gender with body but I wonder if perhaps that is a symptom of the exteriorisation of Catholicism... ... .


\hfill

\end{sffamily}
\end{footnotesize}

